\chapter{Diophantine Equations}

\begin{center}\itshape
$3^2+4^2=5^2$, but $x^4+y^4=z^4$ has no solution, and neither does $y^2=x^3+7$.\\
Four ideas solve a surprising number of equations: factor, reduce modulo $n$, draw a line, descend.
\end{center}

A \emph{Diophantine equation} is a polynomial equation whose integer (or rational) solutions are sought. There is no general method: Hilbert asked for one in 1900, and Matiyasevich proved in 1970 that none can exist. This chapter collects the methods that work surprisingly often.

\section{Pythagorean triples}

A solution of $x^2+y^2=z^2$ in positive integers is \emph{primitive} if $\gcd(x,y,z)=1$. Every solution is a multiple of a primitive one (divide by the gcd), and in a primitive one $x,y$ are coprime (a prime dividing both divides $z^2$, hence $z$).

\begin{think}
The primitive triples with $z<30$ are $(3,4,5)$, $(5,12,13)$, $(15,8,17)$, $(7,24,25)$ and $(21,20,29)$, written with $x$ odd. For each, compute $\frac{z+x}2$ and $\frac{z-x}2$. What do you notice? Can you explain it from $y^2=(z-x)(z+x)$?
\end{think}

\begin{theorem}\label{thm:pyth}
The primitive solutions of $x^2+y^2=z^2$ with $y$ even are exactly
\[
x=m^2-n^2,\qquad y=2mn,\qquad z=m^2+n^2,
\]
with $m>n\ge1$ coprime and of opposite parity.
\end{theorem}
\begin{proof}
These are solutions: $(m^2-n^2)^2+(2mn)^2=(m^2+n^2)^2$. They are primitive: a prime $p$ dividing $x$ and $z$ divides $2m^2$ and $2n^2$; $p\ne2$ since $x$ is odd ($m,n$ of opposite parity), so $p\mid m$ and $p\mid n$, impossible.

Conversely, let $(x,y,z)$ be primitive. $x,y$ are not both odd (else $z^2\equiv2\pmod4$), and not both even, so exactly one is even, say $y$; then $x,z$ are odd. Write $y^2=(z-x)(z+x)$ and $\left(\frac y2\right)^2=\frac{z-x}2\cdot\frac{z+x}2$, a product of two integers. They are coprime: a common divisor divides their sum $z$ and difference $x$. A product of two coprime positive integers is a square only if each is a square (compare the valuations), so $\frac{z+x}2=m^2$ and $\frac{z-x}2=n^2$ with $m>n\ge1$, $\gcd(m,n)=1$. Then $z=m^2+n^2$, $x=m^2-n^2$, $y=2mn$, and $m,n$ have opposite parity because $z$ is odd.
\end{proof}

The key step, ``coprime factors of a square are squares'', is the first example of the \emph{factorisation method}. It will be used again with more exotic factorisations in Chapter~10.

\section{Infinite descent}

\begin{theorem}[Fermat]\label{thm:fermat4}
The equation $x^4+y^4=z^2$ has no solution in positive integers. In particular $x^4+y^4=z^4$ has none.
\end{theorem}
\begin{proof}
Suppose there are solutions, and take one with $z$ minimal. Then $\gcd(x,y)=1$: if a prime $p$ divides $x$ and $y$, then $p^4\mid z^2$, so $p^2\mid z$, and $\left(\frac xp,\frac yp,\frac z{p^2}\right)$ is a smaller solution. So $(x^2,y^2,z)$ is a primitive Pythagorean triple; say $y$ is even. By Theorem~\ref{thm:pyth},
\[
x^2=m^2-n^2,\qquad y^2=2mn,\qquad z=m^2+n^2,
\]
with $m>n\ge1$ coprime of opposite parity. From $x^2+n^2=m^2$ with $x$ odd, $(x,n,m)$ is a primitive Pythagorean triple (any common prime of $n,m$ would divide $\gcd(m,n)=1$), so $n$ is even, $m$ is odd, and $n=2ab$, $m=a^2+b^2$ with $a,b$ coprime. Now $y^2=2mn=4mab$, so $\left(\frac y2\right)^2=mab$, where $m,a,b$ are pairwise coprime ($a,b$ divide $n$, which is prime to $m$). So each of $m,a,b$ is a square: $a=u^2$, $b=v^2$, $m=w^2$. Then
\[
u^4+v^4=a^2+b^2=m=w^2,
\]
a new solution with $w\le w^2=m\le m^2<m^2+n^2=z$. This contradicts the minimality of $z$.
\end{proof}

\begin{remark}
Fermat's Last Theorem, that $x^n+y^n=z^n$ has no solution in positive integers for $n\ge3$, was proved by Wiles (with Taylor) in 1995. Theorem~\ref{thm:fermat4} settles every $n$ divisible by $4$; Euler's case $n=3$ needs the Eisenstein integers of Chapter~10. The general proof is far beyond this series, and we never use it.
\end{remark}

\section{Rational points on conics: the chord method}

\begin{think}
Draw the line through $(-1,0)$ with slope $\frac12$. Where does it meet the circle $x^2+y^2=1$ again? Before computing: why must that second point have rational coordinates? (Think about the sum of the two roots of a quadratic with rational coefficients.)
\end{think}

\begin{proposition}\label{prop:chord}
Let $Q\in\Q[x,y]$ have degree $2$ and no factor of degree $1$ in $\Q[x,y]$, and suppose the conic $Q(x,y)=0$ contains a rational point $P=(x_0,y_0)$. Then every line through $P$ with rational slope $t$ meets the conic in at most one other point, which is rational. Conversely, every rational point other than $P$ is obtained this way, from the line through it and $P$ (or from the vertical line through $P$).
\end{proposition}
\begin{proof}
Substitute $x=x_0+s$, $y=y_0+ts$ in $Q$: we get $g(s)=\alpha s^2+\beta s+\gamma$ with $\alpha,\beta,\gamma\in\Q$, and $\gamma=Q(x_0,y_0)=0$. The polynomial $g$ is not zero: otherwise $Q$ vanishes on the whole line, and writing $\ell=y-y_0-t(x-x_0)$ and $Q=\ell A+Q\big(x,y_0+t(x-x_0)\big)$ (divide by $\ell$ as a polynomial in $y$), the remainder vanishes identically, so $\ell\mid Q$, which we excluded. (Without this hypothesis the statement fails: on $xy=0$ the line $y=0$ through $(0,0)$ lies entirely in the conic.) So $g(s)=s(\alpha s+\beta)$ with $(\alpha,\beta)\ne(0,0)$, whose roots are $0$ (the point $P$) and, if $\alpha\ne0$, $s=-\frac\beta\alpha\in\Q$, giving a rational point. Conversely, if $P'\ne P$ is rational, the line $PP'$ has rational slope (or is vertical), and $P'$ is its second intersection.
\end{proof}

\begin{figure}[ht]
\centering
\fig{ch09-chord-markov}
\caption{Left: the chord method on the unit circle. Each line through $(-1,0)$ with rational slope $t$ meets the circle again at a rational point, and every rational point arises this way; $t=\frac12$ gives the triple $(3,4,5)$, $t=\frac23$ gives $(5,12,13)$. Right: the first levels of the tree of Markov triples (Theorem~\ref{thm:markov}). Each triple $(a,b,c)$ with $c$ largest has two children, $(a,c,3ac-b)$ and $(b,c,3bc-a)$.}
\label{fig:chordmarkov}
\end{figure}

\begin{example}
On $x^2+y^2=1$ with $P=(-1,0)$, the line $y=t(x+1)$ gives $x^2-1+t^2(x+1)^2=0$, i.e.\ $(x+1)\left((x-1)+t^2(x+1)\right)=0$. So the rational points are $(-1,0)$ and
\[
\left(\frac{1-t^2}{1+t^2},\ \frac{2t}{1+t^2}\right),\qquad t\in\Q.
\]
With $t=\frac nm$ and clearing denominators, this recovers Theorem~\ref{thm:pyth}. It is also the substitution $t=\tan\frac\theta2$ of calculus.
\end{example}

A conic without rational points can often be detected by a congruence: $x^2+y^2=3$ has no rational point, because $X^2+Y^2=3Z^2$ has no non-zero integer solution (Exercise~\ref{ex:x2y23z2}). Legendre's theorem (a special case of the Hasse--Minkowski theorem met at the end of Chapter~12 and proved in Tome~12) says that congruence obstructions are the only ones for conics; for cubics this fails (Selmer's $3x^3+4y^3+5z^3=0$).

\section{Factoring with a sum of squares}

\begin{theorem}[Lebesgue, 1869]\label{thm:mordell7}
The equation $y^2=x^3+7$ has no solution in integers.
\end{theorem}
\begin{proof}
If $x$ is even, $y^2\equiv7\pmod8$, impossible since squares are $0,1,4$ modulo $8$. So $x$ is odd. Add $1$ to both sides:
\[
\begin{gathered}y^2+1=x^3+8=(x+2)(x^2-2x+4),\qq x^2-2x+4=(x-1)^2+3.\end{gathered}
\]
As $x-1$ is even, $(x-1)^2+3\equiv3\pmod4$, and it is $\ge3$. A positive integer $\equiv3\pmod4$ has a prime factor $q\equiv3\pmod4$ (a product of primes all $\equiv1\pmod4$ is $\equiv1$, and it is odd). Then $q\mid y^2+1$, so $-1$ is a square modulo $q$, which contradicts $q\equiv3\pmod4$ (Theorem~\ref{thm:eulercrit}).
\end{proof}

The same idea, ``add something to make a sum of two squares, and find a factor $\equiv3\pmod4$'', disposes of many \emph{Mordell equations} $y^2=x^3+k$ (Exercise~\ref{ex:mordell5}). Others, such as $y^2=x^3-2$ (whose only solutions are $(3,\pm5)$), need the arithmetic of $\Z[\sqrt{-2}]$ (Chapter~10).

\section{Vieta jumping}

\begin{think}
$(1,1,1)$ solves $x^2+y^2+z^2=3xyz$. Fix $y=z=1$: then $x$ is a root of $t^2-3t+2$. What is the other root? Now fix the two smallest entries of your new triple and do the same. How far can you go, and does the process ever get stuck?
\end{think}

When an equation is quadratic in one variable, the other root of that quadratic, given by Vi\`ete's relations, is a new solution. This lets us move from solution to solution, and usually to descend.

\begin{theorem}[Markov, 1879]\label{thm:markov}
Every solution of $x^2+y^2+z^2=3xyz$ in positive integers is obtained from $(1,1,1)$ by a finite sequence of permutations and moves $(x,y,z)\mapsto(3yz-x,y,z)$.
\end{theorem}
\begin{proof}
For fixed $y,z$, $x$ is a root of $g(t)=t^2-3yz\,t+(y^2+z^2)$. The other root $x'=3yz-x=\frac{y^2+z^2}x$ is a positive integer, so the move sends solutions to solutions, and it is an involution.

Let $(x,y,z)$ be a solution with $x\ge y\ge z$, other than $(1,1,1)$. We show $x'<x$, so that the move strictly decreases the sum $x+y+z$. Since $z\ge1$, $g(y)=2y^2+z^2-3y^2z\le z^2-y^2\le0$. If $g(y)<0$, then $y$ lies strictly between the roots $x'$ and $x$ of $g$, so $x'<y\le x$. If $g(y)=0$, then $z=1$ and $y=z$, so $y=z=1$ and $x^2-3x+2=0$: $x=2$ (as $x=1$ is excluded), and $x'=1<x$.

So from any solution, repeatedly sorting and applying the move decreases $x+y+z$, which cannot go on forever: we reach a solution where the argument fails, which can only be $(1,1,1)$. Reversing the steps (moves are involutions) builds the original solution from $(1,1,1)$.
\end{proof}

The \emph{Markov numbers} $1$, $2$, $5$, $13$, $29$, $34$, $89$, $169$, $194$, \dots form a tree. The \emph{uniqueness conjecture} (Frobenius, 1913), that the largest element determines the triple, is still open.

\section{Exercises}

\begin{exercise}[$\star$]\label{ex:areaper}
Find all right triangles with integer sides whose area equals their perimeter.
\end{exercise}

\begin{exercise}[$\star$]\label{ex:x2y23z2}
Show that $x^2+y^2=3z^2$ has no solution in integers other than $(0,0,0)$. Deduce that the circle $x^2+y^2=3$ has no rational point.
\end{exercise}

\begin{exercise}[$\star$]\label{ex:catalan23}
Find all $x,y\in\N$ with $2^x-3^y=1$.
\end{exercise}

\begin{exercise}[$\star\star$]\label{ex:mordell5}
Show that $y^2=x^3-5$ has no solution in integers.
\end{exercise}

\begin{exercise}[$\star\star$]\label{ex:eisentri}
Find all rational points of $x^2+xy+y^2=1$, and deduce all triangles with integer sides $a,b,c$ and an angle of $120^\circ$ opposite $c$, that is, $c^2=a^2+ab+b^2$.
\end{exercise}

\begin{exercise}[$\star\star$]\label{ex:markovcoprime}
Show that the three numbers of a Markov triple are pairwise coprime, and list the triples with largest element at most $200$.
\end{exercise}

\begin{exercise}[$\star\star\star$]\label{ex:imo88}
\emph{IMO 1988, Problem 6.} Let $a,b$ be positive integers such that $k=\dfrac{a^2+b^2}{ab+1}$ is an integer. Show that $k$ is a perfect square.
\end{exercise}

\begin{beauty}
Every method of this chapter grows into a theory. Factorisation needs rings where it is unique, or ideals when it is not: Chapters~10 and~11. Congruences become the $p$-adic numbers of Chapter~12, and the principle ``a conic has a rational point if and only if it has one over $\R$ and over every $\Q_p$'' is the Hasse--Minkowski theorem. The chord method on cubics gives the group law on elliptic curves, where Fermat's descent becomes Mordell's theorem: the rational points of an elliptic curve form a finitely generated group (Tome~12).
\end{beauty}

\section{Hints}

\begin{hint}{ex:areaper}
Substitute $c=\frac{ab}2-a-b$ into $a^2+b^2=c^2$ and simplify. The result is linear in each of $a$ and $b$, so it can be written as $(a-\alpha)(b-\alpha)=\beta$ for suitable integers.
\end{hint}
\begin{hint}{ex:x2y23z2}
Work modulo $3$: squares are $0$ or $1$. Show that $x$ and $y$ are multiples of $3$, then $z$ is too, and descend.
\end{hint}
\begin{hint}{ex:catalan23}
For $y\ge1$, reduce modulo $3$ to see that $x$ is even. Then $3^y=(2^{x/2}-1)(2^{x/2}+1)$: two powers of $3$ that differ by $2$.
\end{hint}
\begin{hint}{ex:mordell5}
Add $4$: $y^2+4=x^3-1=(x-1)(x^2+x+1)$. Check $x$ modulo $4$ first, then find a factor $\equiv3\pmod4$.
\end{hint}
\begin{hint}{ex:eisentri}
Use the rational point $(-1,0)$ and the lines $y=t(x+1)$. For the triangles, divide $c^2=a^2+ab+b^2$ by $c^2$.
\end{hint}
\begin{hint}{ex:markovcoprime}
The move $x\mapsto3yz-x$ does not change $\gcd(x,y,z)$, and neither do permutations. Then grow the tree of Figure~\ref{fig:chordmarkov} until the largest entries exceed $200$.
\end{hint}
\begin{hint}{ex:imo88}
Fix $k$ and take a solution $(a,b)$ with $a\ge b$ and $a+b$ minimal. The other root $a'=kb-a=\frac{b^2-k}a$ of $t^2-kbt+b^2-k=0$ is an integer. Show that $a'\ge0$ if $k$ is not a square, and that $a'<b$; what happens when $a'=0$?
\end{hint}

\section{Solutions}

\begin{solution}{ex:areaper}
\idea Write the condition as two equations, $a^2+b^2=c^2$ and $\frac{ab}2=a+b+c$, and eliminate $c$: substitute $c=\frac{ab}2-a-b$ in the first. After dividing by $ab$ what remains is linear in each of $a,b$, and such equations factor (``Simon's favourite factoring trick''): $ab-4a-4b+8=0$ is $(a-4)(b-4)=8$. Then list the divisors of $8$.
\writeup With legs $a,b$ and hypotenuse $c=\frac{ab}2-a-b$, the equation $c^2=a^2+b^2$ becomes $\frac{a^2b^2}4-a^2b-ab^2+2ab=0$; dividing by $ab$ and multiplying by $4$, $ab-4a-4b+8=0$, that is $(a-4)(b-4)=8$. If both factors are negative, they are among $(-1,-8),(-2,-4),(-4,-2),(-8,-1)$, which give $a\le0$ or $b\le0$. So $\{a-4,b-4\}$ is $\{1,8\}$ or $\{2,4\}$: $(a,b)=(5,12)$ or $(6,8)$, with $c=13$ and $10$. Both work: areas $30$ and $24$, perimeters $30$ and $24$.
\end{solution}

\begin{solution}{ex:x2y23z2}
\idea A homogeneous equation with a non-trivial solution has a primitive one, so look for a prime that forces all variables to be divisible by it. Modulo $3$, $x^2+y^2\equiv0$ forces $x\equiv y\equiv0$ (squares are $0,1$). Then $9\mid3z^2$ gives $3\mid z$: descent.
\writeup Suppose $(x,y,z)\ne(0,0,0)$ is a solution; dividing by $\gcd(x,y,z)$ we may assume it is $1$. Modulo $3$ the squares are $0$ and $1$, and $x^2+y^2\equiv0$ forces $x\equiv y\equiv0\pmod3$. Then $9\mid x^2+y^2=3z^2$, so $3\mid z^2$ and $3\mid z$: $3$ divides all three, a contradiction. A rational point $(\frac XZ,\frac YZ)$ on $x^2+y^2=3$ (common denominator $Z\ge1$) would give $X^2+Y^2=3Z^2$ with $Z\ne0$.
\end{solution}

\begin{solution}{ex:catalan23}
\idea Small cases give $(1,0)$ and $(2,1)$. For $y\ge1$, reduce modulo $3$: $2^x\equiv1$ forces $x$ even, and then $2^x-1=(2^{x/2}-1)(2^{x/2}+1)$ is a difference of squares; two powers of $3$ that differ by $2$ can only be $1$ and $3$.
\writeup If $y=0$, $2^x=2$ and $x=1$. Let $y\ge1$. Modulo $3$, $2^x\equiv1$, and $2^x\equiv(-1)^x$, so $x$ is even, $x=2m$ with $m\ge1$. Then $(2^m-1)(2^m+1)=3^y$, so both factors are powers of $3$ (unique factorisation), and they differ by $2$. The only powers of $3$ differing by $2$ are $1$ and $3$ (if both are $\ge3$, their difference is divisible by $3$). So $2^m=2$, $m=1$, $x=2$, $y=1$. The solutions are $(1,0)$ and $(2,1)$.
\end{solution}

\begin{solution}{ex:mordell5}
\idea Imitate Theorem~\ref{thm:mordell7}: make a sum of two squares appear by adding a square to both sides: $y^2+4=x^3-1=(x-1)(x^2+x+1)$. A prime $q\equiv3\pmod4$ cannot divide $y^2+2^2$ unless it divides both $y$ and $2$. So we need a factor $\equiv3\pmod4$; parity considerations modulo $4$ decide which $x$ can occur.
\writeup Modulo $4$: if $x$ is even, $y^2\equiv-5\equiv3$, impossible. If $x\equiv3\pmod4$, $y^2\equiv27-5\equiv2$, impossible. So $x\equiv1\pmod4$. Then $y^2+4=x^3-1=(x-1)(x^2+x+1)$ and $x^2+x+1\equiv1+1+1=3\pmod4$, with $x^2+x+1\ge3$ (as $x^2+x+1>0$ always, and it is $\equiv3$). So $x^2+x+1$ has a prime factor $q\equiv3\pmod4$, and $q\mid y^2+2^2$. If $q\nmid y$, then $\left(2y^{-1}\right)^2\equiv-1\pmod q$, impossible; so $q\mid y$, then $q\mid4$, impossible for odd $q$.
\end{solution}

\begin{solution}{ex:eisentri}
\idea The conic has the obvious rational point $(-1,0)$, so apply the chord method: intersect with $y=t(x+1)$, factor out $x+1$, and solve the remaining linear equation. For triangles, $c^2=a^2+ab+b^2$ means $(\frac ac,\frac bc)$ is a positive rational point on the conic; write $t=\frac nm$ and clear denominators, as for Pythagorean triples.
\writeup With $y=t(x+1)$:
\[\begin{gathered}x^2-1+tx(x+1)+t^2(x+1)^2=0,\qq\text{i.e.}\quad(x+1)\left(x-1+tx+t^2(x+1)\right)=0.\end{gathered}\]
Besides $(-1,0)$, the points are
\[
x=\frac{1-t^2}{1+t+t^2},\qquad y=t(x+1)=\frac{2t+t^2}{1+t+t^2},\qquad t\in\Q,
\]
and the vertical line $x=-1$ adds $(-1,1)$. For $a,b,c>0$ with $c^2=a^2+ab+b^2$, the point $(\frac ac,\frac bc)$ has positive coordinates, which requires $0<t<1$; with $t=\frac nm$, $0<n<m$, the triangle is proportional to
\[
a=m^2-n^2,\qquad b=2mn+n^2,\qquad c=m^2+mn+n^2.
\]
For example $(m,n)=(2,1)$ gives $(3,5,7)$: $9+15+25=49$. Every such triangle is a rational multiple of one of these.
\end{solution}

\begin{solution}{ex:markovcoprime}
\idea Look for a quantity preserved by the move $x\mapsto3yz-x$: the gcd of the triple is one, since $x'\equiv-x$ modulo any common divisor of $y,z$. As $(1,1,1)$ has gcd $1$, every triple does (Theorem~\ref{thm:markov}). Then a prime dividing two of the numbers divides the square of the third, by the equation. For the list, walk down the tree.
\writeup If $d\mid y$ and $d\mid z$, then $d\mid x$ iff $d\mid x'=3yz-x$; so $\gcd(x,y,z)=\gcd(x',y,z)$, and permutations do not change it either. By Theorem~\ref{thm:markov}, every triple has the gcd of $(1,1,1)$, namely $1$. If a prime $p$ divides $x$ and $y$, then $p^2\mid3xyz-x^2-y^2=z^2$, so $p\mid z$: contradiction. Starting from $(1,1,1)$ and applying moves to the largest-but-new element: $(1,1,2)$, $(1,2,5)$, $(1,5,13)$, $(2,5,29)$, $(1,13,34)$, $(1,34,89)$, $(2,29,169)$, $(5,13,194)$; the next ones exceed $200$ (for instance $(1,89,233)$, $(5,29,433)$). Every triple with maximum $\le200$ appears, because by the proof of Theorem~\ref{thm:markov} each triple other than $(1,1,1)$ has a unique parent with a smaller maximum.
\end{solution}

\begin{solution}{ex:imo88}
\idea For fixed $b$ and $k$, the equation $a^2-kb\,a+b^2-k=0$ is quadratic in $a$, so Vi\`ete gives a second solution $a'=kb-a=\frac{b^2-k}a$ with the same $k$. If $k$ is not a square, this jump should produce a smaller solution, and a minimal counterexample gives a contradiction. The places where the descent could stop are $a'<0$ (impossible by a sign check), $a'=0$ (which makes $k=b^2$, a square), and that is exactly what we want.
\writeup Suppose $k$ is not a square, and among all pairs of non-negative integers $(a,b)$ with $\frac{a^2+b^2}{ab+1}=k$, choose one with $a+b$ minimal, and $a\ge b$. If $b=0$ then $k=a^2$, a square; so $b\ge1$. The number $a'=kb-a=\frac{b^2-k}a$ is an integer and satisfies $a'^2-kb\,a'+b^2-k=0$. If $a'<0$, then $a'^2-kba'+b^2-k\ge1+kb+b^2-k>0$, impossible; so $a'\ge0$. If $a'=0$, then $b^2=k$, a square, excluded. So $a'\ge1$ and $(a',b)$ is again a pair with quotient $k$; and $a'=\frac{b^2-k}a<\frac{b^2}a\le a$. Then $a'+b<a+b$ contradicts minimality. Hence $k$ is a square. (Example: $(a,b)=(8,30)$ gives $k=\frac{964}{241}=4$.)
\end{solution}

\section*{Sources and further reading}

\begin{itemize}
\item L.~J.~Mordell, \emph{Diophantine Equations}, Academic Press, 1969. The classic catalogue of methods, including Mordell equations $y^2=x^3+k$.
\item T.~Andreescu, D.~Andrica and I.~Cucurezeanu, \emph{An Introduction to Diophantine Equations}, Birkh\"auser, 2010. Elementary methods with many solved problems, Vieta jumping included.
\item M.~Aigner, \emph{Markov's Theorem and 100 Years of the Uniqueness Conjecture}, Springer, 2013. Everything about Markov triples and the open conjecture.
\item J.~H.~Silverman and J.~Tate, \emph{Rational Points on Elliptic Curves}, 2nd ed., Springer, 2015, Chapter~1. The chord method on cubics and the group law, the next step after this chapter.
\item Y.~V.~Matiyasevich, \emph{Hilbert's Tenth Problem}, MIT Press, 1993. Why no algorithm can decide every Diophantine equation: a meeting point of number theory and computability.
\end{itemize}
